\documentclass{article}

\usepackage[utf8]{inputenc}
\usepackage[T1]{fontenc}
\usepackage{helvet}
\usepackage{geometry}
\usepackage{xcolor}
\usepackage{eso-pic}
\usepackage{fancyhdr}
\usepackage{parskip}
\usepackage{microtype}
\usepackage[english, ukrainian]{babel}
\usepackage{url}
\usepackage{amsmath}
\usepackage{amssymb}
\usepackage{amsthm}
\usepackage{longtable}
\usepackage{booktabs}
\usepackage{hyperref}
\usepackage{titlesec}
\usepackage{tocloft}

\renewcommand{\cftsecfont}{\small\bfseries}
\renewcommand{\cftsubsecfont}{\footnotesize}
\renewcommand{\cftsubsubsecfont}{\scriptsize}
\renewcommand{\cftsecpagefont}{\small\bfseries}
\renewcommand{\cftsubsecpagefont}{\footnotesize}
\renewcommand{\cftsubsubsecpagefont}{\scriptsize}

\definecolor{nato}{HTML}{293757}
\definecolor{footergray}{gray}{0.65}

\geometry{a4paper,left=3cm,right=3cm,top=3cm,bottom=3cm}
\renewcommand{\familydefault}{\sfdefault}
\setcounter{secnumdepth}{3}

\titleformat{\section}
  {\color{nato}\Huge\bfseries\raggedright}{\thesection.}{1em}{}
\titlespacing*{\section}{0pt}{30pt}{12pt}

\titleformat{\subsection}
  {\color{nato}\LARGE\bfseries\raggedright}{\thesubsection.}{1em}{}
\titlespacing*{\subsection}{0pt}{20pt}{8pt}

\titleformat{\subsubsection}
  {\color{nato}\Large\bfseries\raggedright}{\thesubsubsection.}{1em}{}
\titlespacing*{\subsubsection}{0pt}{10pt}{6pt}

\fancyhf{}
\fancyhead[R]{\small\color{footergray} 2 червня 2026}
\fancyfoot[L]{\small\color{footergray} Комплексна система захисту інформації. Цільовий профіль безпеки органів та установ в системі правосуддя (2).}
\fancyfoot[R]{\small\color{footergray}\thepage}

\newcommand\HUGE[1]{\fontsize{40}{40}\selectfont #1}

\renewcommand{\headrulewidth}{0pt}
\renewcommand{\footrulewidth}{0.4pt}
\renewcommand{\footrule}{\color{footergray}\hrule width \textwidth height 0.4pt}

\begin{document}

\pagestyle{fancy}

\begin{titlepage}
  \AddToShipoutPictureBG*{\AtPageLowerLeft{\color{nato}\rule{\paperwidth}{\paperheight}}}
  \color{white}\thispagestyle{empty}\vspace*{4cm}\centering
  {\Huge  \textbf{Комплексна система захисту інформації} \par} \vspace{2cm}
  {\HUGE  \textbf{Цільовий профіль безпеки органів та установ в системі правосуддя (2)} \par} \vspace{2.5cm}
  \vfill
  {\large 2026 \copyright\ Критпографічні Телесистеми \par}
\end{titlepage}

\newpage

\begin{center}
  {\Huge\color{nato}\bfseries Зміст\par}
\end{center}
\vspace{40pt}
\tableofcontents

\begin{abstract}
Цільовий профіль безпеки (ЦПБ) — індивідуальний для конкретної системи підприємства. Саме на його основі створюється/модернізується КЗЗІ. Включає лише додаткові вимоги (відмінності) відносно Галузевого профілю.
\end{abstract}

\section{AC}
Клас заходів захисту AC --- УПРАВЛІННЯ ДОСТУПОМ

Цей клас зосереджується на обмеженні доступу до інформаційних систем, ресурсів та функцій лише для авторизованих користувачів, програм і пристроїв.

\textbf{Перелік заходів захисту:} Розмежування обов'язків (AC-5).

\subsection{РОЗМЕЖУВАННЯ ОБОВ'ЯЗКІВ (AC-5)}
a. Розмежувати і документувати [Призначення: визначені організацією обов'язки окремих осіб].\\ 
b. Установити правила авторизації доступу для підтримки розмежування обов'язків.

\vspace{10pt}
{\footnotesize\bfseries
No: 1\\
Name: ac\_5\_odp\\
Type: string\\
Default: nil
}

{\footnotesize Визначено обов'язки осіб, які потребують розмежування; AC-05[a] обов'язки осіб визначені та задокументовані; AC-05[b] визначено права авторизації доступу до системи для підтримки розмежування обов'язків}




\section{SC}
Клас заходів захисту SC --- ЗАХИСТ ІНФОРМАЦІЙНОЇ СИСТЕМИ ТА

Цей клас охоплює механізми захисту інформації під час її передачі мережами зв'язку, криптографічний захист та ізоляцію критичних компонентів.

\textbf{Перелік заходів захисту:} Id-spe-sc-8-2.

\subsubsection{id-spe-sc-8-2}


\vspace{10pt}
{\footnotesize\bfseries
No: 1\\
Name: sc\_8\_2\_01\\
Type: string\\
Default: nil
}

{\footnotesize КОНФІДЕНЦІЙНІСТЬ ТА ЦІЛІСНІСТЬ ПЕРЕДАЧІ - ПОПЕРЕДНЯ І ПОСТОБРОБКА МЕТА ОЦІНКИ: Визначити, чи:}


{\footnotesize\bfseries
No: 2\\
Name: sc\_8\_2\_odp\\
Type: string\\
Default: nil
}

{\footnotesize Вибрано одне або декілька з наступних ЗНАЧЕНЬ ПАРАМЕТРІВ: \{конфіденційність; цілісність\}}




\end{document}
